\documentclass[border=10pt,tikz]{standalone}
\usepackage{fontspec}
\setmainfont{Apple SD Gothic Neo}
\setsansfont{Apple SD Gothic Neo}
\usepackage{tikz}
\usetikzlibrary{positioning,arrows.meta,shapes,fit,backgrounds,calc}
\begin{document}
% _design.tex — shared TikZ design system for all blog diagrams
% Inputs nothing on its own — meant to be \input{} from each diagram .tex
% AFTER \begin{document}, BEFORE \begin{tikzpicture}.
%
% Provides a unified visual language across GoF / DSA / EMC++ / Beautiful series:
%   - Colors: muted, accessible, color-coded by role
%   - Node styles: consistent sizing, rounded, subtle borders
%   - Edge styles: UML-style inheritance, aggregation, plus dashed alternates
%
% Usage:
%   \documentclass[border=10pt,tikz]{standalone}
%   \usepackage{tikz}
%   \usetikzlibrary{positioning,arrows.meta,shapes,fit,backgrounds}
%   \begin{document}
%   \input{../_design.tex}     % adjust relative path
%   \begin{tikzpicture}[blog]  % apply 'blog' style
%   ...

\definecolor{absfill}{RGB}{220,232,247}   % cool blue
\definecolor{absbord}{RGB}{72,118,182}
\definecolor{confill}{RGB}{249,222,222}   % warm red
\definecolor{conbord}{RGB}{197,93,93}
\definecolor{impfill}{RGB}{223,238,222}   % soft green
\definecolor{impbord}{RGB}{99,154,93}
\definecolor{clifill}{RGB}{251,239,205}   % pale yellow
\definecolor{clibord}{RGB}{200,167,72}
\definecolor{hifill} {RGB}{253,224,200}   % accent orange
\definecolor{hibord} {RGB}{210,135,68}
\definecolor{nodefill}{RGB}{250,250,250}
\definecolor{nodebord}{RGB}{120,120,120}
\definecolor{edgec}  {RGB}{80,80,80}
\definecolor{edgesc} {RGB}{145,145,145}
\definecolor{groupbord}{RGB}{200,200,200}

\tikzset{
  blog/.style={
    font=\sffamily\small,
    every node/.style={font=\sffamily\small}
  },
  % --- node roles ---
  cls/.style={
    rectangle, rounded corners=2pt,
    draw=nodebord, line width=0.4pt,
    minimum height=8mm, minimum width=28mm,
    inner sep=3pt, fill=nodefill, align=center
  },
  abs/.style={cls, draw=absbord, fill=absfill, font=\sffamily\itshape\small},
  con/.style={cls, draw=conbord, fill=confill},
  imp/.style={cls, draw=impbord, fill=impfill},
  cli/.style={cls, draw=clibord, fill=clifill},
  hi/.style={cls, draw=hibord, fill=hifill, font=\sffamily\bfseries\small},
  % --- circle (for trees / graph nodes) ---
  vx/.style={
    circle, draw=absbord, line width=0.4pt,
    minimum size=8mm, fill=absfill, inner sep=0pt
  },
  vxc/.style={vx, draw=conbord, fill=confill},
  vxh/.style={vx, draw=hibord, fill=hifill, font=\sffamily\bfseries},
  % --- decision diamond ---
  q/.style={
    diamond, draw=clibord, line width=0.4pt, fill=clifill,
    inner sep=2pt, aspect=2.4, align=center
  },
  % --- edges ---
  inh/.style={                                  % UML inheritance
    -{Triangle[length=3mm,width=3mm,open]},
    draw=edgec, line width=0.5pt
  },
  agg/.style={                                  % UML aggregation
    {Diamond[length=2.5mm,width=2.5mm,open]}-{Stealth[length=2.2mm]},
    draw=edgec, line width=0.45pt
  },
  uses/.style={                                 % directed reference
    -{Stealth[length=2.2mm]}, draw=edgec, line width=0.5pt
  },
  arr/.style={                                  % plain arrow (graphs/trees)
    -{Stealth[length=2mm]}, draw=edgec, line width=0.45pt
  },
  edge/.style={                                 % undirected (trees)
    draw=edgec, line width=0.45pt
  },
  alt/.style={                                  % dashed: similar/alt
    -{Stealth[length=2mm]}, dashed, draw=edgesc, line width=0.45pt
  },
  ret/.style={                                  % dashed creates/returns
    -{Stealth[length=2mm]}, dashed, draw=edgesc, line width=0.45pt
  },
  thr/.style={                                  % threading link (red dashed)
    ->, dashed, draw=conbord, line width=0.5pt
  },
  % --- group / region ---
  group/.style={
    rectangle, rounded corners=4pt,
    draw=groupbord, line width=0.4pt,
    inner sep=10pt, fill=black!2
  },
  glabel/.style={font=\sffamily\bfseries\small, anchor=north west, inner sep=2pt}
}

% Korean font convenience macro — included by every Hangul diagram
\providecommand{\KOR}{}

% Sources (verified):
%   Linux Documentation/arch/arm64/booting.rst  spin-table: naturally-aligned 64-bit zero-initialised
%       cpu-release-addr inside a /memreserve/ region; CPUs poll it outside the kernel (wfe optional),
%       primary issues sev, jump on non-zero 64-bit little-endian value.
%       psci: kernel issues CPU_ON; DT has a 'psci' node.
%   Linux arch/arm64/kernel/smp_spin_table.c  writeq_relaxed + sev()
%   U-Boot v2024.04 arch/arm/cpu/armv8/spin_table_v8.S  wfe / ldr / cbz / br loop (one implementation)
%   TF-A lib/psci/psci_main.c psci_cpu_on(target_cpu, entrypoint, context_id)
\begin{tikzpicture}[blog,
  bx/.style={cls, minimum width=66mm, minimum height=11mm, align=center, font=\sffamily\small, inner sep=4pt},
  ttl/.style={font=\sffamily\bfseries\small, text=inkc, align=center},
  lbl/.style={font=\sffamily\footnotesize, text=inkc, fill=white, inner sep=2pt},
  dtl/.style={font=\sffamily\footnotesize, text=inkc, align=left, anchor=north west}
]

\node[ttl] at (-5.8, 1.3) {(a) spin-table};
\node[ttl] at ( 5.8, 1.3) {(b) PSCI};

% ---- spin-table ----
\node[hi,  bx] (lp) at (-5.8, 0.0) {Primary (Linux)\\[2pt]release 위치에 entry 주소 기록 후 SEV};
\node[cli, bx] (lm) at (-5.8,-2.2) {\texttt{cpu-release-addr}\\[2pt]64-bit, 초기값 0, 예약 메모리 안};
\node[abs, bx] (ls) at (-5.8,-4.4) {Secondary (커널 밖에서 대기)\\[2pt]release 위치를 polling (WFE는 선택)};
\node[ok,  bx] (lj) at (-5.8,-6.6) {non-zero 값을 읽으면\\[2pt]그 주소로 이동};

\draw[flow] (lp) -- node[lbl, right=2pt]{64-bit LE 값 기록} (lm);
\draw[ret]  (ls.north) -- node[lbl, right=2pt]{읽기} (lm.south);
\draw[flow] (ls) -- (lj);
\draw[selfloop] (ls.east) -- ++(0.5,0) |- node[lbl, pos=0.25, right]{0이면 대기} ($(ls.east)+(0,-0.35)$);

% ---- PSCI ----
\node[hi,  bx] (rp) at (5.8, 0.0) {Primary (Linux)\\[2pt]\texttt{CPU\_ON}(target, entry, context) 호출};
\node[prp, bx] (rc) at (5.8,-2.2) {conduit: SMC 또는 HVC\\[2pt](\texttt{psci} 노드의 \texttt{method})};
\node[con, bx] (rf) at (5.8,-4.4) {PSCI firmware와 플랫폼 코드\\[2pt]대상 CPU 전원·reset 처리};
\node[ok,  bx] (rs) at (5.8,-6.6) {대상 CPU가 요청한 entry에서 시작};

\draw[flow] (rp) -- (rc);
\draw[flow] (rc) -- (rf);
\draw[flow] (rf) -- (rs);

% divider
\draw[draw=mutec, dashed, line width=0.3pt] (0, 1.5) -- (0,-9.2);

% DT description row
\node[dtl] at (-9.1,-7.7) {DT: \texttt{enable-method = "spin-table"}\\[2pt]\phantom{DT: }+ CPU별 \texttt{cpu-release-addr}};
\node[dtl] at ( 2.5,-7.7) {DT: \texttt{enable-method = "psci"}\\[2pt]\phantom{DT: }+ \texttt{psci} 노드};

\end{tikzpicture}
\end{document}
